\documentclass[sigconf,screen,pbalance]{acmart}

\usepackage{booktabs}
\usepackage{microtype}
\setlength{\emergencystretch}{1em}

\setcopyright{cc}
\setcctype{by}
\acmDOI{10.1145/3843282.3844421}
\acmYear{2026}
\copyrightyear{2026}
\acmISBN{979-8-4007-2985-0/2026/10}
\acmConference[AgenticDev '26]{Proceedings of the 1st International Workshop on Agentic AI for Next-Generation Software Development}{October 12--16, 2026}{Munich, Germany}
\acmBooktitle{Proceedings of the 1st International Workshop on Agentic AI for Next-Generation Software Development (AgenticDev '26), October 12--16, 2026, Munich, Germany}
\acmSubmissionID{asews26agenticdevmain-p12-p}
\received{2026-07-10}
\received[accepted]{2026-08-20}

\begin{document}

\shrinkLastPage{75pt}
\balancePageNum{2}

\title{Evaluating Low-Shot Procedural Skill Transfer in Language-Agent
Debugging}

\author{Jiachen Sun}
\correspondingauthor
\orcid{0009-0004-8354-3728}
\affiliation{%
  \institution{ShanghaiTech University}
  \city{Shanghai}
  \country{China}
}
\email{13673527519sjc@gmail.com}

\begin{abstract}
We study whether a natural-language \texttt{SKILL.md} induced from three
debugging trajectories transfers to held-out language-agent debugging tasks.
Our protocol records repair success, diagnosis-edit-test cycles, tokens,
failed patches, and memory-attributed negative transfer. In a single-execution
pilot on six verified PyBugHive \texttt{black} tasks, Auto SKILL.md ties a
generic checklist and Reflexion at 5/6 solved. Task-level comparisons show
lower process cost than reflection baselines but do not establish a stable
advantage over the checklist. One negative-transfer case and an outlier-sensitive
structural control illustrate why memory evaluation needs uncertainty,
process, and misapplication measures rather than solve rate alone.
\end{abstract}

\begin{CCSXML}
<ccs2012>
   <concept>
       <concept_id>10011007.10011074.10011099.10011102.10011103</concept_id>
       <concept_desc>Software and its engineering~Software testing and debugging</concept_desc>
       <concept_significance>500</concept_significance>
   </concept>
   <concept>
       <concept_id>10010147.10010178.10010219.10010221</concept_id>
       <concept_desc>Computing methodologies~Intelligent agents</concept_desc>
       <concept_significance>300</concept_significance>
   </concept>
</ccs2012>
\end{CCSXML}

\ccsdesc[500]{Software and its engineering~Software testing and debugging}
\ccsdesc[300]{Computing methodologies~Intelligent agents}

\keywords{language agents, software debugging, procedural memory, skill
induction, negative transfer}

\maketitle

\section{Motivation and Protocol}
Language-agent memory ranges from verbal reflection to distilled guidelines,
workflows, and procedural repositories
\cite{shinn2023reflexion,zhao2024expel,wang2025workflowmemory,fang2026memp}.
In debugging, however, experience can accelerate a related repair or misdirect
the agent toward an adjacent invariant. We therefore evaluate not only whether
a memory solves a task, but also how it changes repair cost and whether the
trajectory supports memory-attributed negative transfer. This complements
recent work that reports steps or token cost and observes harm from noisy
memory \cite{ouyang2026reasoningbank}.

We induce one fixed artifact from $K=3$ redacted training trajectories. The
induction contract forbids task-specific names and patches and requires three
parts: applicability triggers, an ordered debugging procedure, and failure
modes. The artifact is deployed unchanged on held-out tasks. Five primary
conditions share the same model, task prompts, isolated workspaces, tests,
temperature, and budgets: Auto SKILL.md, a human-written generic checklist,
Reflexion, length-matched Reflexion, and no memory. The checklist already
contains reproduction, localization, minimal editing, and broad validation;
Auto SKILL.md additionally encodes applicability exclusions, expected-versus-
actual invariant compression, and failed-rerun classification.

We evaluate six reproducible PyBugHive \texttt{black} formatter bugs
\cite{antal2024pybughive}. Each task/method pair has one execution. Metrics are
computed from validated trajectories and test reruns. Negative transfer follows
a predefined taxonomy but was assigned by one annotator. For task-level
uncertainty, we use Wilson intervals for solve proportions, paired
task-resampling bootstrap intervals on common-solved tasks, medians, and
leave-one-task-out sensitivity. These do not estimate run-to-run variance.

\begin{table*}[t]
\centering
\caption{Primary same-model pilot. Tokens are solve-conditional means. NT is
the number of memory-attributed negative-transfer flags.}
\label{tab:poster_primary}
\begin{tabular}{lrrrr@{\hspace{1.2em}}r}
\toprule
Method & Solved & Cycles/solved & Tokens/solved & Failed/run & NT \\
\midrule
Auto SKILL.md & 5/6 & 2.00 & 56,867 & 0.83 & 1 \\
Generic checklist & 5/6 & 2.00 & 71,694 & 1.33 & 0 \\
Reflexion memory & 5/6 & 2.60 & 65,915 & 1.33 & 0 \\
Length-matched Reflexion & 4/6 & 2.75 & 75,231 & 1.50 & 0 \\
No memory & 4/6 & 3.00 & 92,031 & 1.17 & 0 \\
\bottomrule
\end{tabular}
\end{table*}

\section{Results and Robustness}
Table~\ref{tab:poster_primary} shows no solve-rate dominance: Auto SKILL.md,
the checklist, and Reflexion each solve 5/6. Wilson 95\% intervals are
0.44--0.97 for 5/6 and 0.30--0.90 for 4/6. On four common-solved tasks,
Auto-minus-baseline mean token differences are $-22.8$k for Reflexion,
$-27.3$k for length-matched Reflexion, and $-11.8$k for the checklist.
Exploratory task-bootstrap intervals are $-44.4$ to $-1.3$k, $-41.4$ to
$-13.1$k, and $-30.3$ to 7.1k, respectively. Cycle differences are $-1.00$,
$-1.00$, and 0.00; the checklist interval ($-0.75$ to 0.75) again shows no
separation. Thus the pilot differentiates Auto SKILL.md more clearly from
reflection memories than from competent human-written advice.

The Auto SKILL.md token median is 52.6k, compared with 56.0k for the checklist
and 68.5k for Reflexion; its leave-one-solved-task-out mean spans 48.0--62.7k.
One Auto run on \texttt{black\_132} is labeled negative transfer: the first
skill-guided patch increased full-suite failures from four to six. This label
is useful diagnostic evidence, not a population-rate estimate.

A supplementary \texttt{gpt-5.5} structure control preserves skill content
but shuffles its organization. Both versions solve 6/6; Auto SKILL.md uses
fewer cycles (1.67 vs. 2.17) and failed patches (0.50 vs. 0.83). Mean tokens
favor shuffled content (102.0k vs. 112.0k), but medians are nearly identical
(86.5k vs. 86.4k), and excluding one outlier reverses the mean ordering. The
control therefore suggests process differences without showing that the
canonical structure is necessary for solvability.

\section{Interpretation and Threats}
The induced artifact is inspectable: relative to the checklist, it makes
applicability boundaries, invariant compression, and failed-rerun categories
explicit. Yet these differences are not equivalent to demonstrated new
knowledge. The tie in solve rate and cycles, together with a checklist token
interval that crosses zero, leaves open whether trajectory induction adds more
than competent human-written procedure. Content-matched section ablations are
needed to identify causal components.

The evidence also has three direct validity limits. First, six selected tasks
and one execution per condition cannot characterize stochastic agent behavior;
our intervals describe task variation only. Second, all primary tasks come
from one formatter project, so within-family transfer may not extend to
repository-wide debugging benchmarks such as SWE-bench
\cite{jimenez2024swebench,yang2024sweagent}. Third, cycle records compress
internal search, token means are outlier-sensitive, and one annotator assigned
negative transfer. Repeated seeds, direct tool instrumentation, and independent
annotation are prerequisites for stronger performance claims.

\section{Conclusion}
The main contribution is an inspectable protocol for evaluating low-shot
procedural-memory transfer, including process cost and misapplication. The
pilot provides initial process evidence but does not establish stable
superiority. Code, artifacts, and validated trajectories are available at
\url{https://github.com/damogu123/trace2skill}.

\begin{figure}[b]
\centering
\begin{minipage}{.96\columnwidth}
\centering
\textbf{3 redacted debugging trajectories}\\[2pt]
$\Downarrow$ trajectory-to-skill induction\\[2pt]
\textbf{SKILL.md}: triggers $\mid$ ordered procedure $\mid$ failure modes\\[2pt]
$\Downarrow$ fixed deployment\\[2pt]
\textbf{6 held-out bugs} under one harness\\
Controls: checklist $\mid$ Reflexion $\mid$ length-matched $\mid$ no memory\\
Measures: solve $\mid$ cycles/tokens $\mid$ failed patches $\mid$ negative transfer
\end{minipage}
\Description{Three redacted debugging trajectories are compressed into a
SKILL.md containing triggers, an ordered procedure, and failure modes. The
fixed artifact and four controls are evaluated on six held-out bugs using
success, efficiency, failed-patch, and negative-transfer metrics.}
\caption{Low-shot procedural-memory induction and held-out evaluation.}
\label{fig:poster_protocol}
\end{figure}

\bibliographystyle{ACM-Reference-Format}
\bibliography{references}

\end{document}
